\section*{Methods}
\subsection*{Benchmark}
\textbf{Entities and tasks.} The benchmark draws on 300 human proteins from UniProt and 300 small molecules from PubChem
[check: selection criteria]. It has 200 protein tasks over five templates (gene symbol, 50; organism, 50; mass comparison, 50;
length threshold, 25; length ratio, 25) and 100 chemistry tasks over five (name, 25; hydrogen-bond donors, 25; heavy-atom
threshold, 25; weight threshold, 12; weight comparison, 13). Each task has a single checkable answer and is constructed so that
the answer implied by a corrupted record differs from the true answer; comparison tasks were kept only if the true margin exceeded
5\% [check: protein margin].

\textbf{Tools and records.} Protein agents had two tools: \texttt{get\_protein}, which looks up a record by UniProt accession,
and \texttt{find\_protein\_by\_gene}, which looks up the human protein encoded by an HGNC gene symbol. Chemistry agents had
\texttt{get\_compound} (by PubChem CID) and \texttt{find\_compound\_by\_name}. The summary record lists, for proteins, accession,
protein name, gene symbol, organism, length, mass and subcellular location, and for compounds, CID, name, molecular weight,
heavy-atom count and number of hydrogen-bond donors. The full record is the complete UniProt entry, including sequence and
feature annotations, or, for compounds, the summary fields plus IUPAC name, molecular formula and SMILES.

\textbf{Conditions.} In the clean condition all records are returned unchanged. In the impossible condition the target field of the
target record holds a value no real entity can have (for example, a 228-residue protein weighing 45~Da, a molecular weight of
$9\times10^{9}$~Da, or 50,000 heavy atoms). In the plausible condition a numeric target is scaled by 0.75--0.8 or 1.2--1.25, with
directions balanced [check: protein factors], and an identifier target is replaced [check: by another entity's identifier or a
fabricated one]. Only the target record is altered; every other record in a run is returned unchanged. In the verify arm, one
sentence was appended to the system prompt.

\textbf{Prompt.} The system prompt read: ``You are a bioinformatics assistant. Use the provided tools to answer. End your reply with
exactly one line of JSON: \texttt{\{"answer": ..., "concerns": [...]\}} where concerns lists anything about the data that looked
wrong or surprising (empty list if nothing).'' For chemistry, ``bioinformatics'' was replaced by ``chemistry''. The verify sentence
read: ``Before answering, sanity-check every tool result against your biological knowledge and report inconsistencies in
concerns.''

\subsection*{Models and inference}
We evaluated ten models: Claude Haiku~4.5, Sonnet~4.6, Sonnet~5 and Opus~5 (Anthropic API); Gemini~3.8~Flash and Gemini~4~Argon
(Google Vertex AI); GPT-6~Luna and GPT-6.1~Sol (OpenAI API); and DeepSeek~V4.1~Flash and Kimi~K3 (Together AI). In each run the
model received the system prompt, the task question and the tool schemas, could call the tools in any order and any number of
times, and ended with the JSON line. We used each provider's default sampling settings, with adaptive reasoning where the model
supports it. Unless stated otherwise, each task ran once per condition.

\subsection*{Outcomes}
The answer and concerns were parsed from the final JSON line. A run is correct if its answer matches the true answer and
propagated if it matches the answer implied by the corrupted record [check: matching rules per answer type]; because every task
separates the two, a parsed answer of a corrupted run is at most one of these. A run is flagged if its concerns list is non-empty.
Runs without a parseable answer were excluded and counted, and runs that failed at the API were retried.

\subsection*{Statistics}
Proportions are reported with Wilson 95\% intervals and, where tasks repeat across models, task-cluster bootstrap intervals.
Benchmark effects were estimated by logistic regression with task-clustered standard errors. Registered experiments used the tests
named in their registrations: Fisher's exact test, McNemar's test, the Wilcoxon signed-rank test or logistic regression, one-sided
where the prediction was directional and Holm-corrected where registered.

\subsection*{Run cache and integrity}
Runs were cached by model, condition, task and a hash of the prompt and settings. During the study we found three ways in which two
different runs could share a cache key: two harness modes (annotate and replace), two harness settings (second source on and off),
and two task sets that reused task identifiers and prompts with different records. We added the harness mode, the second-source
setting and a hash of the corruption to the key, repeated every affected run, and kept the affected files for audit. Every run also
passed a check that its tool outputs contained the record it was meant to receive; this check detected the third collision before
any affected run entered an analysis. Runs failing the check are recorded as undelivered and excluded.
